\documentclass[9pt]{extarticle}
\usepackage[margin=0.5in, headheight=14pt]{geometry}
\usepackage{enumitem}
\usepackage{booktabs}
\usepackage{hyperref}
\usepackage{xcolor}
\usepackage{fancyhdr}
\usepackage{titlesec}

\hypersetup{colorlinks=true, linkcolor=blue, urlcolor=blue}
\setlength{\parindent}{0pt}
\setlength{\parskip}{3pt}
\titlespacing*{\section}{0pt}{8pt}{3pt}
\titlespacing*{\subsection}{1.5em}{6pt}{2pt}
\setlist[itemize]{leftmargin=1.5em}
\setlist[enumerate]{leftmargin=1.5em}

\title{\vspace{-1.2cm}\huge\textbf{A3: Infrastructure and Dependencies}\vspace{-0.5em}}
\author{\normalsize Autobook}
\date{\vspace{-1.5em}}

\pagestyle{fancy}
\fancyhf{}
\fancyhead[L]{\small A3: Infrastructure and Dependencies}
\fancyhead[R]{\small Autobook}
\fancyfoot[C]{\small \thepage}
\renewcommand{\headrulewidth}{0.4pt}
\renewcommand{\footrulewidth}{0pt}

\begin{document}
\maketitle
\thispagestyle{fancy}

What tools we use for A3 and why each was chosen.

\begin{itemize}[nosep]
  \item Nanochat fork --- training codebase
  \item Modal --- GPU compute
  \item Weights \& Biases --- experiment tracking
  \item YAML configs --- run specification
\end{itemize}

\leftskip=0pt
\section*{Nanochat Fork --- Training Code}
\leftskip=1.5em

We fork Karpathy's nanochat (Feb 2026) so that P2 can modify the architecture without touching the original code.

\textbf{URL}: \url{https://github.com/seonghyunban/nanochat}

\subsection*{Branches}
\leftskip=3em

\begin{itemize}[nosep]
  \item \texttt{master} --- original nanochat, frozen. All baseline runs use this. Do not push here.
  \item \texttt{p2} --- P2's architecture changes, branched from \texttt{master}.
\end{itemize}

\subsection*{Tag}
\leftskip=3em

\begin{itemize}[nosep]
  \item \texttt{baseline-v0} --- immutable snapshot of \texttt{master}. Safety net: even if someone pushes to \texttt{master}, this tag never moves.
\end{itemize}

\subsection*{How Modal Uses the Fork}
\leftskip=3em

The container clones the fork once at image build and caches it. Before every run, it does \texttt{git fetch} to pull new commits, then checks out the branch/tag the config asks for. No container rebuild needed when P2 pushes new code.

\leftskip=0pt
\section*{Modal --- GPU Compute}
\leftskip=1.5em

A3 needs single-GPU jobs. Modal is the simplest platform for this: write a Python function, pick a GPU, run it. \$500 credits ready via promo code, per-second billing, and Volumes to save checkpoints between runs.

\subsection*{Why Not Others}
\leftskip=3em

Anyscale (\$1{,}000) is built around Ray for distributed workloads --- more setup than needed for single-GPU jobs, kept as P4 backup. AWS (\$500) has the most overhead for simple jobs (EC2, SSH, drivers, storage). Snowflake (\$400 trial) is a data warehouse, not a training platform.

\subsection*{Which GPU}
\leftskip=3em

Nanochat requires BF16, so A10G is the minimum. A100/H100 are faster and often cheaper in total cost due to shorter wall time.

\begin{tabular}{@{}llrr@{}}
\toprule
Part & GPU & Est.\ Time & Est.\ Cost \\
\midrule
P2 (4 ablation runs) & A100 & 1--5 hrs & \$3--13 \\
P3 (2 training stages + eval) & A100 & 1--3 hrs & \$3--8 \\
P4 (full nanochat) & H100 & 6--20 hrs & \$24--79 \\
\bottomrule
\end{tabular}

\leftskip=0pt
\section*{Weights \& Biases --- Experiment Tracking}
\leftskip=1.5em

Dashboard for training curves, loss values, and run configs. Already built into nanochat --- no extra code needed.

\begin{itemize}[nosep]
  \item \textbf{Project}: \texttt{490-autobook-a3} (shared across all parts)
  \item \textbf{Tags}: each run tagged \texttt{p2}, \texttt{p3}, or \texttt{p4}
  \item \textbf{Git hash}: every run logs which nanochat commit it used
\end{itemize}

\leftskip=0pt
\section*{YAML Configs --- Run Specification}
\leftskip=1.5em

Each experiment is a YAML file whose fields map directly to nanochat's command-line flags (\texttt{--depth}, \texttt{--max-seq-len}, etc.). No extra config framework needed --- nanochat already parses its own arguments.

\begin{itemize}[nosep]
  \item \texttt{nanochat\_ref} field picks which branch or tag to run
  \item Comments in the YAML explain each field --- no separate docs needed
\end{itemize}

\end{document}
